\section{FENOMENI ONDULATORI}

\subsection{Introduzione}

In questo capitolo discuteremo alcuni dei pi\`u importanti fenomeni
caratteristici delle grandezze che obbediscono all'equazione delle onde: si
tratta di quei fenomeni (interferenza, diffrazione \dots) che storicamente hanno
permesso di stabilire la natura ondulatoria della luce e, successivamente, dei
raggi X, ma che essendo conseguenze dell'equazione delle onde si manifestano per
molti altri sistemi, come il suono, le onde in un liquido, ecc.. Dato che la
nostra discussione riguarda solo i fenomeni elettromagnetici, gli argomenti di
questo capitolo rientrano in quello pi\`u vasto dell'Ottica ondulatoria (o
Ottica fisica), un capitolo della fisica di grande interesse ma anche molto
esteso: basti solo dire che il suo naturale completamento \`e l'Ottica
geometrica.\\
Il nostro obiettivo non \`e quello di una discussione ragionevolmente esauriente
di questi argomenti (sia pure con le limitazioni che le tecniche a nostra
disposizione ci impongono), ma una discussione mirata alla comprensione degli
argomenti che normalmente vengono utilizzati come introduzione ai fondamenti
della Meccanica Quantistica: questo spiega sia la trattazione molto semplificata
del problema dell'interferenza che quella, praticamente solo a livello
qualitativo, del problema della diffrazione. Un altro argomento che possiamo
portare per giustificare le nostre scelte, \`e che lo strumento matematico
``naturale'' per la discussione di questi problemi, in particolare quello della
diffrazione, non \`e tanto quello delle formule approssimate di
Huygens--Fresnel, ma quello della trasformata di Fourier, un argomento che
solitamente viene trattato nei corsi degli anni superiori al secondo.

\subsection{L'interferenza}

L'interferenza, come altri fenomeni ondulatori, trae origine dalla linearit\`a
delle equazioni di Maxwell nel vuoto, per cui ogni combinazione lineare di
soluzioni \`e anch'essa una soluzione (principio di sovrapposizione).\\
Sia data un'onda piana monocromatica che si propaga in direzione dell'asse $x$:
\begin{equation}
  \vec E(x,t)
  =\big(0,\,E_{0y}\cos(kx-\omega t+\phi),\,E_{0z}\cos(kx-\omega t+\psi)\big)
  \tag{13.1}
\end{equation}
Se consideriamo uno schermo piano in direzione ortogonale alla direzione di
propagazione (per esempio una lastra fotografica), l'intensit\`a che incide
sullo schermo, data dalla (12.32), \`e indipendente da $y$ e da $z$ (oltre che
dalla posizione dello schermo): lo schermo viene quindi ``illuminato''
uniformemente.\\
Supponiamo ora che sullo schermo incidano due onde piane $\vec E_1(x,t)$ ed
$\vec E_2(x,t)$ con la stessa frequenza e che si propagano entrambe in direzione
dell'asse $x$:
\begin{equation}
  \begin{aligned}
    &\vec E_1(x,t)=\big(0,\,E_{1y}^0\cos(kx-\omega t+\phi_1),
    \,E_{1z}^0\cos(kx-\omega t+\psi_1)\big)\\
    &\vec E_2(x,t)=\big(0,\,E_{2y}^0\cos(kx-\omega t+\phi_2),
    \,E_{2z}^0\cos(kx-\omega t+\psi_2)\big)
  \end{aligned}
  \tag{13.2}
\end{equation}
il campo \ $\vec E(x,t)=\vec E_1(x,t)+\vec E_2(x,t)$ \ \`e ancora un'onda piana
monocromatica, e ancora lo schermo viene illuminato uniformemente. Tuttavia, in
generale, l'intensit\`a con cui lo schermo viene illuminato \big(cio\`e
l'intensit\`a dell'onda $\vec E(x,t)$\big), non \`e la somma delle intensit\`a
delle due onde piane $\vec E_1(x,t)$ ed $\vec E_2(x,t)$: infatti, tenuto conto
che
\begin{equation}
  \frac{1}{T}\int_t^{t+T}\cos(\omega t'-\alpha)\cos(\omega t'-\beta)\,dt'
  =\frac{1}{2}\cos(\alpha-\beta)
  \tag{13.3}
\end{equation}
si ha
\begin{equation}
  \begin{aligned}
    I&\equiv\frac{c}{4\pi}\overline{|\vec E_1(x,t)+\vec E_2(x,t)|^2}
    =\frac{c}{4\pi}\Big(\overline{\vec E_1(x,t)^2}+\overline{\vec E_2(x,t)^2}
    +2\overline{\vec E_1(x,t)\cdot\vec E_2(x,t)}\Big)\\
    &=I_1+I_2+\frac{c}{4\pi}\big(E_{1y}^0E_{2y}^0\cos(\phi_1-\phi_2)
    +E_{1z}^0E_{2z}^0\cos(\psi_1-\psi_2)\big)
  \end{aligned}
  \tag{13.4}
\end{equation}
per cui
\begin{equation}
  I_1+I_2-2\sqrt{I_1I_2}\le I\le I_1+I_2+2\sqrt{I_1I_2}
  \tag{13.5}
\end{equation}
Fissati $I_1$ ed $I_2$, il massimo ed il minimo per $I$ si hanno quando
\ $\cos(\phi_1-\phi_2)=\cos(\psi_1-\psi_2)=\pm1$ \ ed i vettori $\vec E_{01}$ ed
$\vec E_{02}$ sono paralleli. Quando $I\ne I_1+I_2$ si dice che le due onde
danno luogo al fenomeno dell'\textbf{interferenza}; se $I<I_1+I_2$ si dice che
l'interferenza \`e distruttiva, mentre nel caso opposto si dice che \`e
costruttiva.\\
L'intensit\`a pu\`o essere nulla (interferenza completamente distruttiva) quando
\ $\vec E_1(x,t)=-\vec E_2(x,t)$ \ (e quindi $\vec E(x,t)=0$), ed \`e massima
(interferenza completamente costruttiva) quando \ $\vec E_1(x,t)=\vec E_2(x,t)$
\ \big(e quindi $\vec E(x,t)=2\vec E_1(x,t)$\big), ed in tal caso \ $I=4I_1$.
Nel primo caso si dice che le due onde sono \textbf{in controfase}, mentre nel
secondo caso si dice che le onde sono \textbf{in fase}.\\
Due onde entrambe polarizzate linearmente in direzioni fra loro ortogonali non
interferiscono:
\begin{equation}
  \begin{aligned}
    &\vec E_1(x,t)=\vec E_{01}\cos(kx-\omega t+\phi_1),
    \quad\ \vec E_2(x,t)=\vec E_{02}\cos(kx-\omega t+\phi_2)\quad\ \Rightarrow\\
    &I=I_1+I_2+\frac{c}{4\pi}\vec E_{01}\cdot\vec E_{02}\cos(\phi_1-\phi_2)
    =I_1+I_2
  \end{aligned}
  \tag{13.6}
\end{equation}
in quanto (per ipotesi) \ $\vec E_{01}\cdot\vec E_{02}=0$.

Supponiamo ora di avere due sorgenti poste nei punti $x=0$, $y=0,\,z=\pm d$, e
lo schermo nel piano $x=D$, con $D\gg d$.\\
Prendiamo per i campi emessi dalle due sorgenti quelli dati dalla (12.46);
limiteremo la regione di osservazione sullo schermo alla zona in cui $y,z\ll D$,
per cui in tutti i punti di detta regione i raggi $\vec r_1$ e $\vec r_2$
provenienti dalle due sorgenti formano fra di loro un piccolo angolo, di modo
che possiamo assumere che sia i campi elettrici che quelli magnetici delle due
onde siano paralleli (onde polarizzate nello stesso modo). Con le ipotesi fatte,
in tutta la regione di osservazione $\theta\approx90^\circ$
\ ($\sen\theta\approx1$) per entrambe le onde provenienti dalle due sorgenti; in
detta regione \ $\vec E_{1,2}\parallel z$ \ e
\fig{fig13_1}{13.1}
\begin{equation}
  \begin{aligned}
    &E_{1z}(\vec r,t)=\frac{A}{r_1}\sen(kr_1-\omega t+\phi)\,,
    \qquad E_{2z}(\vec r,t)=\frac{A}{r_2}\sen(kr_2-\omega t)\\
    &r_{1,2}=\sqrt{D^2+y^2+(z\pm d)^2}
  \end{aligned}
  \tag{13.7}
\end{equation}
Inoltre, nella regione dello schermo alla quale siamo interessati, consideriamo
costanti ed uguali a $D$ le distanze $r_1$ ed $r_2$ dalle sorgenti che compaiono
nei denominatori della (13.7); non possiamo considerarle costanti quando
compaiono nella forma $kr$ come
argomenti del seno, in quanto, per piccole lunghezze d'onda, cio\`e grandi $k$,
anche se $r$ varia poco ($\Delta r\ll r$), $kr$ pu\`o variare di parecchi
multipli di $2\pi$: in particolare vedremo che la distribuzione dell'intensit\`a
sullo schermo dipende proprio dalla differenza fra $kr_1$ e $kr_2$.\\
Allora, posto \ $I_0=(c/8\pi)(A/D)^2$ \ (l'intensit\`a sullo schermo dovuta a
ciascuna sorgente),
\begin{equation}
  \begin{aligned}
    I(y,z)&=2I_0\big(\overline{\sen^2(kr_1-\omega t+\phi)}
    +\overline{\sen^2(kr_2-\omega t)}+\\
    &\qquad+\overline{2\sen(kr_1-\omega t+\phi)\sen(kr_2-\omega t)}\big)=\\
    &=2I_0\big(1+\cos\big(k(r_1-r_2)+\phi\big)\big)
  \end{aligned}
  \tag{13.8}
\end{equation}
L'intensit\`a di illuminazione in un punto dello schermo dipende quindi sia
dallo sfasamento $\phi$ fra le due onde, che dalla differenza $r_1-r_2$ dei
``cammini ottici'': se per esempio $\phi=0$, l'interferenza \`e completamente
costruttiva nei punti in cui i cammini ottici differiscono per un numero intero
di lunghezze d'onda ($\lambda=2\pi/k$), ed \`e completamente distruttiva dove la
differenza fra i cammini ottici \`e un multiplo dispari di semi--lunghezze
d'onda.\\
Siccome nelle nostre ipotesi \ $|r_1-r_2|\ll D$, conviene riscrivere la (13.8)
in una forma pi\`u espressiva, approssimando \ $r_1-r_2$ \ con
\ $2d\,\partial r/\partial z$ \ e $r$ con $D$:
\begin{equation}
  I(y,z)\approx2I_0\Big(1+\cos\big(\frac{2kdz}{D}+\phi\big)\Big)
  \tag{13.9}
\end{equation}
L'intensit\`a di illuminazione sullo schermo varia quindi in direzione dell'asse
$z$ da 0 a $4I_0$ con un alternarsi periodico di \textbf{frange} pi\`u chiare
dove l'interferenza \`e costruttiva \ $(\cos(2kdz/D+\phi)>0)$, e di frange pi\`u
scure dove l'interferenza \`e distruttiva \ $(\cos(2kdz/D+\phi)<0)$. La
larghezza $\Delta z$ di ogni frangia, cio\`e la distanza fra due massimi o due
minimi di illuminazione, \`e data da
\begin{equation}
  \frac{2kd\Delta z}{D}=2\pi
  \quad\ \Rightarrow\quad\ \Delta z=\frac{\pi D}{kd}=\frac{\lambda D}{2d}
  \tag{13.10}
\end{equation}
ed i centri delle frange chiare si trovano nei punti dati da
\begin{equation}
  \frac{2kdz_n}{D}+\phi=2\pi n
  \quad\ \Rightarrow\quad\ z_n=\big(n-\frac{\phi}{2\pi}\big)\frac{\lambda D}{2d}
  \qquad n=\pm1,\pm2,\cdots
  \tag{13.11}
\end{equation}
per cui se viene variato lo sfasamento $\phi$ fra le due onde, tutta la figura
di interferenza (cio\`e l'insieme delle frange) si sposta rigidamente.

\subsection{La diffrazione}

Supponiamo che un'onda piana (p.es.\ monocromatica) incida ortogonalmente su uno
schermo nel quale \`e praticato un foro di forma arbitraria: circolare,
rettangolare, \dots\\
Ci si potrebbe aspettare che l'onda prosegua oltre il foro con un profilo, nel
piano ortogonale alla direzione di propagazione, uguale alla forma del foro, e
quindi che si abbia un passaggio netto dalla zona illuminata alla zona d'ombra.
Ci\`o non \`e possibile, in quanto $\vec E$ e $\vec B$, pur in assenza di
sorgenti, subirebbero una discontinuit\`a nel passaggio da una zona all'altra.
Si potrebbe obiettare che una separazione netta fra zona illuminata e zona
d'ombra non \`e fisica, in quanto il bordo del foro (anche se ``otticamente ben
lavorato'') non \`e una linea geometrica, ma una regione fisica in cui si passa
dal vuoto
allo schermo su distanze dell'ordine (almeno) delle dimensioni atomiche, e
quindi si potrebbe supporre che l'onda oltre il foro, relativa a qualsiasi
componente del campo elettrico o magnetico, sia della forma
\begin{equation}
  F(x,y,z,t)=AS(y,z)\cos(kx-\omega t)\qquad\ x>0
  \tag{13.12}
\end{equation}
dove \ $A\cos(kx-\omega t)$ \ \`e l'onda incidente sullo schermo (posto nel
piano $x=0$), $S(y,z)$ \`e la funzione caratteristica del foro opportunamente
regolarizzata, cio\`e una funzione derivabile quanto basta, che vale 1 quando
$y,z$ sono interni al foro, e nulla al di fuori di una regione appena superiore
a quella del foro.\\
Neppure la (13.12) pu\`o essere soluzione dell'equazione delle onde (12.4),
indipendentemente dal fatto che $F$ sia una componente del campo
elettromagnetico, o una qualsiasi grandezza che deve obbedire all'equazione
delle onde (onda sonora, onda in un liquido, \dots); infatti
\begin{equation}
  \Delta\,F-\frac{1}{c^2}\frac{\partial^2F}{\partial t^2}
  =A\Big(\frac{\partial^2S}{\partial y^2}+\frac{\partial^2S}{\partial z^2}\Big)
  \cos(kx-\omega t)
  \tag{13.13}
\end{equation}
quindi affinch\'e $F$ sia soluzione dell'equazione delle onde \`e necessario che
\begin{equation}
  \frac{\partial^2S}{\partial y^2}+\frac{\partial^2S}{\partial z^2}=0
  \tag{13.14}
\end{equation}
\`E facile vedere che la (13.14) ammette solo la soluzione identicamente nulla:
basta, per esempio, pensare $S$ come un potenziale elettrostatico in tre
dimensioni (costante rispetto a $x$), privo di sorgenti.\\
Ci\`o significa che la forma della soluzione oltre il foro \`e molto complicata,
e non \`e nemmeno possibile determinarla esplicitamente (occorre soddisfare le
condizioni al contorno imposte dalla presenza dello schermo).\\
Non possiamo andare molto a fondo nella discussione di questo problema: anche le
trattazioni approssimate richiedono una discussione non banale per essere
giustificate.\\
La cosa pi\`u semplice che possiamo dire, e che \`e alla base delle trattazioni
approssimate (principio di Huygens -- Fresnel dedotto dalle formule di
Kirchhoff), \`e che il foro si comporta come una sorgente estesa di onde
elettromagnetiche: \`e come se ogni elemento infinitesimo del foro si
comportasse come un'antenna (parallela al campo elettrico incidente) ed essa,
mantenuta in oscillazione dal campo incidente ed in fase con esso, emettesse
un'onda del tipo della (12.46); il campo oltre il foro \`e quindi la
sovrapposizione dei campi emessi da questo insieme infinito di antenne
infinitesime, e l'interferenza fra tutte queste onde determina la distribuzione
dell'intensit\`a.\\
Il risultato \`e che, per esempio nel caso di un foro circolare di raggio
$d\gg\lambda$, l'intensit\`a \`e praticamente diversa da zero in un cono di
apertura \ $\theta\approx\lambda/d$, attorno alla direzione dell'onda incidente
(il fattore numerico di proporzionalit\`a fra $\theta$ e $\lambda/d$ dipende
dalla forma del foro): questo \`e il fenomeno della \textbf{diffrazione},
evidentemente tanto pi\`u rilevante quanto pi\`u piccolo \`e il foro, e tanto
maggiore la lunghezza d'onda. Per angoli maggiori di $\theta$, l'intensit\`a
tende a zero con oscillazioni di ampiezza de-
\fig{fig13_2}{13.2}
crescente (figura 13.2). Quindi il passaggio dalla zona illuminata alla zona
d'ombra \`e graduale e non monotono, per cui su uno schermo in un piano
$x=D\gg d$, l'immagine del foro appare come una macchia illuminata di dimensioni
$\approx\theta D$, contornata da un insieme di anelli scuri e anelli chiari di
intensit\`a rapidamente decrescente (figura di diffrazione del foro): se
$\theta=\lambda/d$ non \`e trascurabile, le dimensioni della figura di
diffrazione hanno poco a che fare con le dimensioni del foro.\\
Se il foro \`e di forma rettangolare di lati $a,b$; \ $a<b$, la figura di
diffrazione risulta allungata nella direzione del lato corto
\ ($\Delta y\approx\lambda D/a$), e accorciata nella direzione del lato lungo
\ ($\Delta z\approx\lambda D/b$).

\paragraph{Esempio}
\begin{quote}
  Calcoliamo la dimensione $\Delta$ della frangia centrale prodotta da un
  forellino di 0.1\,mm, attraversato da luce di lunghezza d'onda
  $\lambda=6000$\,\AA{} (luce gialla), su uno schermo a distanza $D=2$\,m dal
  foro:
  \begin{equation}
    \Delta\approx\frac{6\times10^{-5}\cdot2\times10^2}{10^{-2}}
    \simeq1\,\mathrm{cm}\gg0.1\,\mathrm{mm}
    \tag{13.15}
  \end{equation}
\end{quote}
Se l'onda (piana e monocromatica) incide ortogonalmente su uno schermo sul quale
sono praticati due fori di dimensioni piccole rispetto alla loro distanza $2d$,
\underline{qualitativamente} la situazione \`e identica a quella discussa alla
fine del paragrafo precedente: il campo oltre lo schermo \`e la sovrapposizione
dei campi dati dalla (13.7) (con $\phi=0$), cio\`e dei campi di due antenne nei
punti $z=\pm d$ (ma ci\`o \`e vero solo se \`e valida l'approssimazione di
dipolo, cio\`e se la dimensione dei fori \`e molto minore della lunghezza
d'onda), ed il risultato \`e (sui piani a distanza $D\gg d$) una figura di
interferenza simile a quella data dalla (13.9): se la condizione che il foro sia
piccolo rispetto alla lunghezza d'onda non \`e soddisfatta (p.es.\ nel caso
della luce), l'intensit\`a delle frange decresce allontanandosi dalla frangia
centrale. Si noti che la larghezza delle frange, data dalla (13.10), coincide
con la dimensione della frangia centrale della figura di diffrazione di un foro
di dimensione $2d$: in effetti tutta la figura di interferenza pu\`o essere
considerata come la figura di diffrazione prodotta da una fenditura di altezza
$2d$ in direzione dell'asse $z$; il fatto che questa fenditura sia in buona
parte ostruita si riflette sull'intensit\`a, e in misura molto minore sulla
posizione, delle frange adiacenti a quella centrale.

\subsection{La luce}

Il processo di emissione della luce da parte degli atomi \`e un processo
quantistico, che non pu\`o essere spiegato con le leggi dell'elettromagnetismo
classico: in un atomo gli elettroni hanno una accelerazione quindi, come ogni
carica accelerata, dovrebbero emettere onde elettromagnetiche perdendo in tal
modo energia fino a fermarsi; siccome l'unica configurazione di equilibrio \`e
quella in cui gli elettroni sono appiccicati al nucleo (non esistenza di
posizioni di equilibrio stabile nel vuoto), ne segue che tutti gli elettroni
finirebbero per cadere sul nucleo, e quindi le dimensioni stabili dell'atomo
sarebbero quelle nucleari, cio\`e circa $10^{-12}\div10^{-13}$\,cm (``collasso
per irraggiamento'').\\
Uno dei risultati fondamentali della Meccanica Quantistica \`e che l'energia di
un atomo (energia cinetica pi\`u energia potenziale) non pu\`o essere
arbitraria, ma pu\`o solo avere determinati valori discreti $E_n$ limitati
inferiormente: \ $E_0<E_1<E_2<\cdots$. Queste
energie permesse $E_n$ sono i \textbf{livelli energetici} dell'atomo ed il
livello di energia pi\`u bassa $E_0$ \`e detto \textbf{livello} (o
\textbf{stato}) \textbf{fondamentale}.\\
La luce viene emessa quando l'atomo passa da un livello energetico $E_n$ ad un
altro di energia inferiore $E_m$: la frequenza della luce emessa nella
transizione $E_n\to E_m$ \`e legata al ``salto quantico''
\ $\Delta E\equiv E_n-E_m$ \ (che, per la conservazione dell'energia, \`e anche
l'energia del ``quanto'' (\textbf{fotone}) di luce emessa), dalla relazione di
Einstein
\begin{equation}
  E_\mathrm{fotone}=\Delta E=h\nu
  \qquad\qquad h=6.6\times10^{-27}\,\mathrm{erg\,s}
  \tag{13.16}
\end{equation}
$h$ \`e una costante fondamentale tipica dei processi quantistici, nota come
costante di Planck.\\
Quando l'atomo si trova nello stato fondamentale non pu\`o pi\`u irraggiare,
quindi lo stato fondamentale \`e lo stato di equilibrio di un atomo.\\
Le transizioni fra due livelli avvengono in un tempo dell'ordine di
$10^{-8}\div10^{-9}$\,s, per cui l'estensione spaziale dell'onda emessa \`e $c$
volte il suddetto tempo: $30\div300$\,cm; un'onda di estensione spaziale (e
quindi temporale) limitata \`e detta un \textbf{pacchetto d'onde} (la
terminologia nasce dalla teoria della trasformata di Fourier).

La luce emessa da una sorgente luminosa monocromatica, per esempio da una
lampada al sodio (luce gialla), non \`e quindi un'unica onda monocromatica, ma
\`e un insieme statistico di fotoni (in quanto ogni atomo emette indipendente
dagli altri), cio\`e un insieme statistico di pacchetti d'onda con la stessa
frequenza, ma con polarizzazioni diverse (esistono diversi stati dell'atomo che
corrispondono allo stesso livello energetico, e la polarizzazione del fotone
emesso dipende dallo stato di partenza e da quello di arrivo nella transizione
fra due livelli energetici). Quindi un fascio di luce naturale non \`e mai
polarizzato.\\
Esistono in natura, ma soprattutto in commercio, sostanze plastiche (polaroid)
che permettono di ottenere luce polarizzata linearmente: queste sostanze hanno
la propriet\`a di essere completamente trasparenti alla luce polarizzata
linearmente in una certa direzione (detta ``asse ottico''), di assorbire
completamente la luce polarizzata in direzione ortogonale all'asse ottico, e di
avere un comportamento \textbf{lineare}: il campo in uscita dipende linearmente
dal campo in entrata (quello che stiamo descrivendo \`e un polaroid ideale: la
trasparenza e l'assorbimento non sono mai completi).\\
Consideriamo un'onda piana che si propaga in direzione dell'asse $x$, e
disponiamo ortogonalmente alla direzione di propagazione un foglio polaroid;
prendiamo l'asse $z$ (per esempio) nella direzione dell'asse ottico. L'onda
piana \`e descritta dalla (13.1) e pu\`o essere considerata come la
sovrapposizione di un'onda polarizzata in direzione dell'asse $y$:
\ $\vec E_1(x,t)=\big(0,\,E_{0y}\cos(kx-\omega t+\phi),\,0\big)$, e di un'onda
polarizzata in direzione dell'asse $z$:
\ $\vec E_2(x,t)=\big(0,\,0,\,E_{0z}\cos(kx-\omega t+\psi)\big)$. Grazie alla
linearit\`a del polaroid, la prima viene completamente assorbita, e la seconda
completamente trasmessa, quindi l'onda uscente dal polaroid \`e polarizzata in
direzione dell'asse $z$, cio\`e dell'asse ottico del polaroid.\\
Se indichiamo con $I$ \ l'intensit\`a dell'onda incidente e con $I_\mathrm{tr}$
quella dell'onda trasmessa, e con $\vec E_0$ il vettore $(0,\,E_{0y},\,E_{0z})$,
si ha
\begin{equation}
  I=\frac{c}{8\pi}\vec E_0^2\,,\quad\ I_\mathrm{tr}=\frac{c}{8\pi}E_{0z}^2\,;
  \quad\ \Rightarrow\quad\ I_\mathrm{tr}=I\cos^2\alpha
  \tag{13.17}
\end{equation}
dove $\alpha$ \`e l'angolo fra il vettore $\vec E_0$ e l'asse ottico. La (13.17)
\`e nota come \textbf{Legge di Malus}.\\
Se sul polaroid incidono diverse onde piane (o diversi pacchetti d'onda) con
polarizzazioni diverse, le onde trasmesse sono tutte polarizzate in direzione
dell'asse ottico, ed in tal modo, a scapito dell'intensit\`a, si pu\`o ottenere
un fascio di luce polarizzata: ci\`o pu\`o essere dimostrato sperimentalmente
disponendo due polaroid con assi ottici paralleli e verificando che tutta la
luce trasmessa dal primo attraversa anche il secondo (oppure, se gli assi ottici
sono ortogonali, che il sistema dei due polaroid \`e completamente opaco).

Nasce per\`o un problema: se il fascio incidente \`e un fascio di fotoni, cosa
accade al singolo fotone, cio\`e al singolo pacchetto d'onde, che non sia
polarizzato n\'e parallelamente n\'e ortogonalmente all'asse ottico del polaroid
ma, per esempio, sia polarizzato linearmente in una direzione che forma un
angolo $\alpha$ con l'asse ottico? Secondo quanto sopra, ne dovrebbe essere
trasmessa una frazione, cio\`e l'energia del fotone trasmesso per la legge di
Malus dovrebbe essere uguale a $h\nu\cos^2\alpha$, e quindi, se vale la
relazione di Einstein (13.16) che lega l'energia di un fotone con la sua
frequenza, il fotone trasmesso dovrebbe avere una frequenza
\ $\nu'=\nu\cos^2\alpha$. Che le cose non stiano cos\`i ce lo dice l'ovvia
constatazione che se sul polaroid incide luce gialla (o rossa, o verde), esce
luce gialla (o rossa, o verde), cio\`e luce con la stessa frequenza. Questo \`e
un esempio tipico del conflitto fra la natura ondulatoria della luce, espressa
dal principio di sovrapposizione, e la natura corpuscolare introdotta dalla
relazione (13.16) di Einstein: il fotone ha un'energia, una quantit\`a di moto
$E/c$ \ e soprattutto \`e, come una particella, un'entit\`a
\underline{indivisibile}; il modo in cui la Meccanica Quantistica risolve questo
problema costituisce uno dei pi\`u profondi cambiamenti dei paradigmi su cui si
basa la fisica classica.

\subsection{L'interferometro di Mach e Zehnder}

In tempi recenti sono stati realizzati esperimenti di interferenza nei quali
anzich\'e un fascio di luce veniva utilizzato un fascio di particelle, pi\`u
precisamente neutroni. Questi esperimenti forniscono una prova molto convincente
della natura ondulatoria delle particelle e quindi delle previsioni della
Meccanica Quantistica.\\
L'interferometro utilizzato negli esperimenti con i neutroni \`e molto simile
all'interferometro di Mach e Zehnder, utilizzato per la luce: diamo quindi una
breve descrizione di questo interferometro.
\fig{fig13_3}{13.3}
L'interferometro di Mach e Zehnder, rappresentato in figura 13.3, consiste di
quattro specchi fra loro paralleli e disposti ai vertici di un rettangolo e
formano con i lati di questo rettangolo angoli di 45$^\circ$. Due di questi
specchi ($s_2$ ed $s_3$) sono normali specchi, mentre gli altri due ($s_1$ ed
$s_4$) sono specchi semi--trasparenti: un fascio di luce incidente su di essi
viene parzialmente riflesso e parzialmente trasmesso (``beam--splitters'');
supporremo che il 50\% dell'intensit\`a venga riflessa e il 50\% trasmessa.\\
Un fascio di luce, provenendo da sinistra, incide su $s_1$: il fascio trasmesso
si riflette su $s_3$ e incide su $s_4$, dal quale viene parzialmente riflesso
verso il rivelatore $R_1$ e parzialmente trasmesso verso il rivelatore $R_2$; il
fascio riflesso da $s_1$ si riflette su $s_2$ ed infine
viene da $s_4$ parzialmente trasmesso verso $R_1$, e parzialmente riflesso verso
$R_2$. Quindi, su ciascuno dei due rivelatori giunge sia parte del fascio
proveniente dal ramo $s_2-s_4$, sia parte del fascio proveniente dal ramo
$s_3-s_4$, e questi due fasci interferiscono fra di loro: le intensit\`a $I_1$ e
$I_2$ registrate dai due rivelatori $R_1$ ed $R_2$ sono il risultato
dell'interferenza fra i due fasci che giungono su ciascuno di essi: se uno
(qualsiasi) dei due fasci venisse intercettato, p.es.\ mediante l'interposizione
di un materiale assorbente in uno dei due cammini, i due rivelatori $R_1$ ed
$R_2$ registrerebbero la stessa intensit\`a, 1/4 dell'intensit\`a del fascio
incidente.\\
Prendiamo l'asse $x$ nella direzione del fascio incidente; supponiamo che la
luce incidente sia un'onda piana, che per semplicit\`a supponiamo linearmente
polarizzata (ma questa ipotesi non \`e affatto necessaria), in modo che possiamo
scrivere l'onda come
\begin{equation}
  F(x,t)=F_0\cos(kx-\omega t)
  \tag{13.18}
\end{equation}
dove $F$ \`e l'unica componente del campo elettrico diversa da zero.\\
Siccome l'intensit\`a dell'onda \`e quadratica nell'ampiezza $F_0$, l'ampiezza
dell'onda trasmessa da uno specchio semi--trasparente e quella dell'onda
riflessa sono entrambe ridotte di un fattore $\sqrt2$ rispetto all'ampiezza
dell'onda incidente. Quindi l'onda che arriva su $R_1$ lungo il percorso
$s_1\to s_2\to s_4$, avendo incontrato entrambi gli specchi semi--trasparenti,
\`e
\begin{equation}
  \frac{1}{2}F_0\cos(kx-\omega t+\phi_1)
  \tag{13.19}
\end{equation}
e quella che arriva su $R_1$ lungo il percorso $s_1\to s_3\to s_4$ \ \`e
\begin{equation}
  \frac{1}{2}F_0\cos(kx-\omega t+\phi_2)
  \tag{13.20}
\end{equation}
Se i cammini ottici \ $s_1\to s_2\to s_4$ \ e \ $s_1\to s_3\to s_4$ \ fossero
rigorosamente identici, $\phi_1=\phi_2$, ma da un lato \`e impossibile calibrare
l'interferometro in modo che la differenza fra i due cammini ottici sia molto
minore della lunghezza d'onda, e d'altro lato pu\`o essere utile introdurre
deliberatamente uno sfasamento fra le due onde, quindi nelle (13.19) e (13.20)
manteniamo $\phi_1\ne\phi_2$ e poniamo $\phi=\phi_1-\phi_2$.\\
Quindi, se \ $I\equiv cF_0^2/8\pi$ \ \`e l'intensit\`a dell'onda incidente su
$s_1$, l'intensit\`a su $R_1$ \`e
\begin{equation}
  \begin{aligned}
    I_1&=\frac{c}{4\pi}\frac{1}{4}F_0^2\,
    \overline{\big(\cos(kx-\omega t+\phi_1)+\cos(kx-\omega t+\phi_2)\big)^2}\\
    &=\frac{1}{2}I\,(1+\cos\phi)
  \end{aligned}
  \tag{13.21}
\end{equation}
Di conseguenza, dato che \ $I_1+I_2=I$, deve essere
\begin{equation}
  I_2=\frac{1}{2}I\,(1-\cos\phi)
  \tag{13.22}
\end{equation}
L'intensit\`a su ciascuno dei rivelatori pu\`o quindi variare da 0 a $I$, al
variare dello sfasamento $\phi$: lo sfasamento pu\`o essere variato con
continuit\`a inserendo in uno dei rami dell'interferometro una sottile lastrina
di vetro, e variandone l'inclinazione, in modo da cambiare lo spessore del vetro
attraversato dal fascio: se $n$ \`e l'indice di rifrazione del vetro, lo
sfasamento introdotto \`e dato da \ $\delta\phi=k(n-1)d=2\pi(n-1)d/\lambda$,
dove $d$ \`e lo spessore attraversato.
